%% kapa: exam=1 round=27 year=2016 period=1 form=A subject=civil_law subject_ko=민법 seq=4 number=4
%% source: Q-Net 감정평가사 기출문제 (cst00309) artlSeq=5169402 file=extracted/1008/5169402/제27회 감정평가사 1차 문제/제27회 감정평가사 1교시문제(A형).pdf
%% conversion: kordoc 4.24.51 -> tools/gen1.py (auto); review=proofread-sonnet; flags=none
\begin{question}{4}
민법상 법인의 기관에 관한 설명으로 옳지 않은 것은?
\begin{choices}
  \item 특별대리인은 임시기관으로 법인의 대표기관이다.
  \item 이사에 의해 선임된 대리인은 법인의 대표기관이 아니다.
  \item 감사는 필요기관으로 그 성명과 주소를 등기하여야 한다.
  \item 이사가 없는 경우에 이로 인하여 손해가 생길 염려가 있는 때에는 법원은 이해관계인이나 검사의 청구에 의해 임시이사를 선임하여야 한다.
  \item 법인의 불법행위가 성립하는 경우, 그 가해행위를 한 이사 기타 대표자는 자기의 손해배상책임을 면하지 못한다.
\end{choices}
\end{question}
\answer{3}
